% Project presentation (10--15 min): DL Complexity.
% Style follows the provided Beamer example: default theme, 16:9, blue headings,
% restrained blocks, frame numbers, and no navigation symbols.
% Build: tectonic dl_complexity_presentation.tex
\documentclass[aspectratio=169,notheorems]{beamer}
\usetheme{default}
\usefonttheme{professionalfonts}
\setbeamersize{text margin left=8mm,text margin right=8mm}
\setbeamertemplate{navigation symbols}{}
\setbeamertemplate{footline}{%
  \leavevmode
  \hbox{\begin{beamercolorbox}[
    wd=\paperwidth,ht=2.5ex,dp=1.2ex,center
  ]{page number in head/foot}%
    \insertframenumber
  \end{beamercolorbox}}}
\setbeamertemplate{itemize item}{--}
\setbeamertemplate{itemize subitem}{--}

\usepackage{fontspec}
\setsansfont{Arial}
\setmainfont{Arial}
\XeTeXlinebreaklocale "en"
\usepackage{amsmath,amssymb,mathtools,booktabs,array}
\usepackage{microtype}
\newcolumntype{L}[1]{>{\raggedright\arraybackslash}p{#1}}

\definecolor{myblue}{RGB}{25,75,130}
\setbeamercolor{structure}{fg=myblue}
\setbeamercolor{frametitle}{fg=myblue}
\setbeamerfont{frametitle}{size=\large,series=\bfseries}
\setbeamertemplate{blocks}[rounded][shadow=false]
\setbeamercolor{block title}{fg=myblue,bg=myblue!12}
\setbeamercolor{block body}{bg=myblue!4}

\newcommand{\E}{\mathbb{E}}
\newcommand{\KL}{\mathrm{KL}}
\newcommand{\NLL}{\mathrm{NLL}}
\newcommand{\code}[1]{\texttt{#1}}
\newcommand{\take}[1]{%
  \par\vspace{0.35em}{\color{myblue}\hrule height 0.6pt}%
  \vspace{0.25em}\noindent\textbf{#1}\par}

\title{DL Complexity}
\subtitle{Theoretically informed deep learning model complexity estimation}
\author{Alyabushev \and Dementiev \and Kurdyukov \and Chirkov}
\date{Technical Meeting 1}

\begin{document}

% ============================================================ TITLE
\begin{frame}[plain]
\centering
\vspace*{1.2em}
{\LARGE\bfseries\color{myblue} DL Complexity}\\[0.5em]
{\large Theoretically informed deep learning model complexity estimation}\\[2.0em]
Alyabushev, Dementiev, Kurdyukov, Chirkov\\[0.4em]
{\small Python package: \code{dl\_complexity}}\\[2.0em]
Technical Meeting 1
\end{frame}

% ============================================================ PROBLEM
\begin{frame}{Problem statement}
\small
Parameter count alone neither determines the complexity of the represented function nor
explains generalization in overparameterized networks.

\medskip
\textbf{Research question:} which theoretically motivated estimators are most strongly
associated with held-out NLL and the generalization gap across dataset sizes,
architectures, and regularization settings?

\medskip
\begin{itemize}\setlength{\itemsep}{0.45em}
  \item Implement several estimator families behind a common interface.
  \item Compare them under one controlled experimental protocol.
  \item Report the score, computational cost, rank stability, assumptions, and warnings.
\end{itemize}

\take{The project compares useful approximations; it does not assume one universal
true measure of neural-network complexity.}
\end{frame}

% ============================================================ SCOPE
\begin{frame}{Project deliverables and scope}
\small
\begin{columns}[T]
\begin{column}{0.48\textwidth}
\begin{block}{Deliverables}
\begin{itemize}\setlength{\itemsep}{0.3em}
  \item an installable Python library;
  \item a common benchmark runner;
  \item reproducible configurations;
  \item comparison tables and plots;
  \item documentation and a demonstration workflow.
\end{itemize}
\end{block}
\end{column}
\begin{column}{0.48\textwidth}
\begin{block}{First release}
\begin{itemize}\setlength{\itemsep}{0.3em}
  \item supervised classification;
  \item small MLPs and CNNs;
  \item synthetic data, MNIST, and Fashion-MNIST;
  \item PyTorch models;
  \item checkpoint-based evaluation where possible.
\end{itemize}
\end{block}
\end{column}
\end{columns}

\medskip
\textbf{Estimator-specific work:} posterior sampling, curvature estimation, compression,
or repeated training is performed when the method requires it.

\textbf{Out of scope:} large language models, distributed training, automatic model
selection, and proofs of new generalization bounds.

\textbf{Current status:} planning; the estimators have not yet been implemented.
\end{frame}

% ============================================================ INTERFACE
\begin{frame}[shrink=4]{Project name and common interface}
\small
\textbf{Project name:} DL Complexity. \hfill
\textbf{Python package:} \code{dl\_complexity}.

\medskip
\begin{block}{Primary interface}
\code{ComplexityEstimator.estimate(context) -> ComplexityResult}
\end{block}

\renewcommand{\arraystretch}{1.16}
\begin{tabular}{@{}L{0.27\textwidth}L{0.67\textwidth}@{}}
\toprule
\textbf{Type} & \textbf{Purpose} \\
\midrule
\code{EstimationContext} & model, dataset, split, likelihood, training artifacts, seed \\
\code{ComplexityResult} & native value, valid nats/bits conversion, per-sample value,
target, direction \\
 & components, assumptions, warnings, runtime, peak memory, metadata \\
\code{EstimatorConfig} & prior, precision, blocks, posterior samples, compression settings \\
\bottomrule
\end{tabular}

\medskip
With natural logarithms, AIC/BIC/HQIC/WAIC are in deviance units: divide by $2$ for an
NLL-like value in nats and by $2\ln 2$ for bits. This conversion does not make them
literal prefix-code lengths.

\take{One output schema does not imply one semantics: predictive criteria, evidence
approximations, operational codes, and risk bounds remain distinct.}
\end{frame}

% ============================================================ ARCHITECTURE
\begin{frame}{Library architecture, integration, and technology stack}
\small
\centering
\begin{tabular}{c@{$\to$}c@{$\to$}c@{$\to$}c}
\fbox{\parbox{0.18\textwidth}{\centering
model/data\\[-0.1em]
adapters}}
&
\fbox{\parbox{0.18\textwidth}{\centering
\code{Registry}\\[-0.1em]
validation}}
&
\fbox{\parbox{0.20\textwidth}{\centering
\code{Estimator}\\[-0.1em]
four families}}
&
\fbox{\parbox{0.18\textwidth}{\centering
\code{Result}\\[-0.1em]
report}}
\end{tabular}

\medskip
$\downarrow$\\[-0.2em]
\fbox{\parbox{0.74\textwidth}{\centering
\code{BenchmarkRunner}: train/load models $\rightarrow$ estimate $\rightarrow$
aggregate $\rightarrow$ tables and plots}}

\medskip
\begin{columns}[T]
\begin{column}{0.48\textwidth}
\textbf{Integration}
\begin{itemize}\setlength{\itemsep}{0.25em}
  \item \code{torch.nn.Module}, \code{DataLoader};
  \item likelihood and posterior-sample adapters;
  \item optional curvature and compression backends;
  \item JSON/CSV-compatible result records.
\end{itemize}
\end{column}
\begin{column}{0.48\textwidth}
\textbf{Technology stack}
\begin{itemize}\setlength{\itemsep}{0.25em}
  \item Python 3.11+, PyTorch;
  \item NumPy, SciPy, pandas, scikit-learn;
  \item pytest, coverage, ruff, static typing;
  \item MkDocs; Jupyter or Streamlit for the demo.
\end{itemize}
\end{column}
\end{columns}
\end{frame}

% ============================================================ ALGORITHM 1
\begin{frame}[shrink=9]{Algorithms I: information criteria and two-part MDL}
\footnotesize
Let $\widehat L$ be the maximized \emph{total} likelihood, $k$ the number of estimated
free parameters, and $n$ the number of independent observations. Natural logs are used.
\[
\mathrm{AIC}=-2\log\widehat L+2k,\quad
\mathrm{BIC}=-2\log\widehat L+k\log n,\quad
\mathrm{HQIC}=-2\log\widehat L+2k\log\log n.
\]
\textbf{WAIC} is a posterior predictive criterion:
\[
\mathrm{WAIC}=-2(\mathrm{lppd}-p_{\mathrm{WAIC}}),\quad
\mathrm{lppd}=\sum_i\log\E_{\theta\mid D}p(y_i\mid\theta),\quad
p_{\mathrm{WAIC}}=\sum_i\operatorname{Var}_{\theta\mid D}\log p(y_i\mid\theta).
\]
\textbf{WBIC} approximates Bayesian free energy using a tempered posterior:
\[
p_\beta(\theta\mid D)\propto p(D\mid\theta)^\beta p(\theta),\quad
\beta=1/\log n,\quad
\mathrm{WBIC}=\E_\beta[-\log p(D\mid\theta)].
\]
\textbf{Two-part MDL} is an operational code:
\[
L_{\mathrm{2p}}=L(\theta_\delta)+L(D\mid\theta_\delta),
\]
where a prefix code, quantization rule $\delta$, model structure, and hyperparameters
must be specified.

\take{AIC/BIC/HQIC are baselines for deep networks; their regular-model assumptions
usually fail because neural networks are singular and non-identifiable.}
\end{frame}

% ============================================================ ALGORITHM 2
\begin{frame}[shrink=5]{Algorithms II: Bayesian coding and marginal evidence}
\small
\begin{block}{Variational approximation}
For prior $p(\theta)$ and variational approximation $q_\phi(\theta)$,
\[
\mathcal L_{\mathrm{var}}
=\E_{q_\phi}[-\log p(D\mid\theta)]
+\KL\!\left(q_\phi(\theta)\,\|\,p(\theta)\right)
=-\mathrm{ELBO}\geq-\log p(D).
\]
The coding interpretation of the KL term requires an explicit scheme such as bits-back coding.
\end{block}

\begin{block}{KFAC--Laplace approximation}
At the MAP estimate $\widehat\theta$, let
$H=\nabla_\theta^2[-\log p(D,\theta)]_{\widehat\theta}$ be positive definite. Then
\[
-\log p(D)\approx-\log p(D,\widehat\theta)
+\tfrac12\log\det H-\tfrac d2\log(2\pi).
\]
KFAC replaces full curvature by layerwise Kronecker-factored Fisher/GGN blocks.
\end{block}

\textbf{The estimate depends on} the prior, variational family, curvature definition,
damping, and the validity of a single-mode Gaussian approximation.
\end{frame}

% ============================================================ ALGORITHM 3
\begin{frame}[shrink=4]{Algorithms III: prequential coding}
\small
Ordered data are partitioned into blocks $0=n_0<n_1<\dots<n_m=n$. The first labels
use a uniform code; each later block is encoded by a model trained only on previously
decoded labels:
\[
L_{\mathrm{preq}}
=n_1\log_2 K
-\sum_{j=2}^{m}\sum_{i=n_{j-1}+1}^{n_j}
\log_2 p_{\widehat\theta(D_{<j})}(y_i\mid x_i).
\]

\begin{center}
\fbox{uniform labels}
$\longrightarrow$
\fbox{train on $D_1$, encode $D_2$}
$\longrightarrow$
\fbox{train on $D_{1:2}$, encode $D_3$}
$\longrightarrow\cdots$
\end{center}

\begin{itemize}\setlength{\itemsep}{0.34em}
  \item Inputs $x_i$ are shared side information; only labels are encoded.
  \item Architecture, optimizer, ordering, block schedule, calibration protocol, and
        reproducible training procedure must be shared or encoded.
  \item Implement a blockwise baseline and an online/replay variant; report every
        training run and wall-clock cost.
\end{itemize}

\textbf{Validation:} encoder and decoder reproduce the same probabilities without
future-label leakage. Shorter per-example length on readily learnable data is a
diagnostic, not a theorem.
\end{frame}

% ============================================================ ALGORITHM 4
\begin{frame}[shrink=3]{Algorithms IV: compression-based complexity}
\small
A trained network is replaced by a discrete compressed representation $C$ using
pruning, quantization, or low-rank factorization. Two related quantities are kept separate:
\[
\underbrace{L(C)+L(D\mid C)}_{\text{two-part data code}}
\qquad\text{and}\qquad
\underbrace{\ell(C)}_{\text{model description in bits}}.
\]
For a prefix-coded hypothesis and bounded loss, a schematic Occam bound is
\[
R(C)\ \leq\ \widehat R(C)+
\sqrt{\frac{\ell(C)\ln 2+\ln(1/\delta)}{2n}},
\]
with the exact constants determined by the selected theorem.

\begin{itemize}\setlength{\itemsep}{0.32em}
  \item The bound applies to the compressed model; relating it to the original model
        requires a controlled loss or margin distortion.
  \item The encoder, decoder, precision, sparsity indices, and side information are
        fixed by the protocol and counted in the code.
  \item A serialized \code{.pt} file size is not a theoretical description length.
\end{itemize}

\textbf{Validation:} exact round-trip decoding, reproducible bit count, and measured
accuracy, NLL, and margin change after compression.
\end{frame}

% ============================================================ BENCHMARK
\begin{frame}[shrink=4]{Benchmark protocol and hypotheses}
\small
\renewcommand{\arraystretch}{1.13}
\begin{tabular}{@{}L{0.27\textwidth}L{0.67\textwidth}@{}}
\toprule
\textbf{Component} & \textbf{Controlled or varied} \\
\midrule
Data & synthetic classification, MNIST, Fashion-MNIST; fixed train/validation/test splits \\
Models & MLP/CNN; width, depth, regularization, multiple random seeds \\
Training & optimizer, stopping rule, compute budget; MLE-like baseline where required \\
Outputs & native score, justified nats/bits conversion, components, assumptions, warnings \\
Resources & wall-clock time, number of training runs, peak memory \\
\bottomrule
\end{tabular}

\medskip
\textbf{Two empirical hypotheses}
\begin{enumerate}\setlength{\itemsep}{0.2em}
  \item total criterion or code length, normalized per example, versus held-out NLL;
  \item consistently normalized penalty or model-description component versus the generalization gap
        (test NLL minus train NLL, in evaluation mode).
\end{enumerate}

\textbf{Analysis:} Spearman correlation, confidence intervals, rank stability across
seeds and sample sizes, and estimator cost relative to model training. Validation data
select hyperparameters; the test set is used only for final evaluation.

\take{Compare lower-is-better ranks within each estimator target; unit conversion alone
does not create semantic comparability.}
\end{frame}

% ============================================================ POC
\begin{frame}{Proof-of-concept demonstration}
\small
\textbf{PoC for Technical Meeting 2}
\begin{enumerate}\setlength{\itemsep}{0.32em}
  \item Train 3--5 MLPs of different widths on one fixed MNIST split.
  \item Through the common API, compute parameter count, BIC, two-part MDL, and a
        blockwise prequential code.
  \item Save the configuration, seed, checkpoints, score components, time, and memory.
  \item Plot complexity against held-out NLL and the generalization gap.
\end{enumerate}

\medskip
\begin{block}{Basic user workflow}
\code{result = registry.create("bic", config).estimate(context)}\\
\code{report = BenchmarkRunner(experiment).run(estimators)}
\end{block}

\textbf{Final demo output:} a comparable result table with native units, rank order,
score decomposition, assumptions, warnings, and runtime. One CLI or notebook workflow
must reproduce the result from a clean environment.

\take{The PoC validates the interface and the end-to-end data path; it is not yet the
final empirical study.}
\end{frame}

% ============================================================ ROLES
\begin{frame}[shrink=5]{Team responsibilities and algorithm ownership}
\scriptsize
\renewcommand{\arraystretch}{1.20}
\begin{tabular}{@{}L{0.13\textwidth}L{0.28\textwidth}L{0.50\textwidth}@{}}
\toprule
\textbf{Member} & \textbf{Algorithm family} & \textbf{Additional responsibilities and deliverable} \\
\midrule
Alyabushev & AIC, BIC, HQIC, WAIC, WBIC, two-part MDL &
Planning; benchmarks and PoC; final comparison tables. \\
Dementiev & Variational coding; KFAC--Laplace evidence &
Planning; library packaging and common API; tests and coverage above 90\%. \\
Kurdyukov & Blockwise and online/replay prequential coding &
Planning; public API documentation; reproducible final demonstration. \\
Chirkov & Compression-based score and generalization bound &
Planning; blog post; three-to-five-page technical report. \\
\bottomrule
\end{tabular}

\medskip
\textbf{Cross-review:}
Alyabushev $\rightarrow$ Dementiev $\rightarrow$ Kurdyukov $\rightarrow$ Chirkov
$\rightarrow$ Alyabushev.

\medskip
Each algorithm owner supplies the mathematical definition and assumptions, interface
implementation, minimal example, unit tests, and documentation section.
\end{frame}

% ============================================================ MILESTONES
\begin{frame}[shrink=4]{Milestones and acceptance criteria}
\small
\renewcommand{\arraystretch}{1.15}
\begin{tabular}{@{}L{0.21\textwidth}L{0.70\textwidth}@{}}
\toprule
\textbf{Stage} & \textbf{Verifiable deliverable} \\
\midrule
Technical Meeting 1 & scope, definitions, units, architecture, stack, API, roles, PoC plan \\
Technical Meeting 2 & end-to-end PoC; draft documentation, blog post, and technical report \\
Checkpoint & installable package; consistent code style; workflows; cross-review \\
Technical Meeting 3 & all estimator families; benchmark; coverage above 90\%; demo; final report \\
\bottomrule
\end{tabular}

\medskip
\begin{block}{Definition of done}
\begin{itemize}\setlength{\itemsep}{0.23em}
  \item installation and execution succeed in a clean environment;
  \item every estimator implements one contract and returns a structured result;
  \item benchmark configurations fix data, training, units, and evaluation protocol;
  \item tests cover formulas, edge cases, and end-to-end integration;
  \item documentation states assumptions and provides reproducible examples;
  \item the report separates established theory, implementation choices, and empirical findings.
\end{itemize}
\end{block}
\end{frame}

% ============================================================ RISKS
\begin{frame}[shrink=4]{Main risks and controls}
\small
\renewcommand{\arraystretch}{1.16}
\begin{tabular}{@{}L{0.43\textwidth}L{0.50\textwidth}@{}}
\toprule
\textbf{Risk} & \textbf{Control} \\
\midrule
Regular-model criteria are applied to singular neural networks &
treat them as baselines; state violated assumptions and use a controlled MLE-like run \\
WAIC/WBIC require posterior inference &
use a dedicated adapter; do not report a score when required draws are absent \\
Laplace is sensitive to prior, curvature, and damping &
save the decomposition and run a sensitivity analysis \\
Prequential coding requires many repeated fits &
start with small models; fixed block schedule; online/replay variant \\
Compression results depend on the exact code and distortion &
fix encoder/decoder; count side information; verify round-trip and quality change \\
Scores use different targets and native scales &
retain target and native value; convert only when justified; compare stratified ranks \\
\bottomrule
\end{tabular}

\medskip
\textbf{Resource constraint:} first validate one minimal estimator per family on an MLP;
add CNNs and the full experiment grid only after the shared interface is stable.
\end{frame}

% ============================================================ SUMMARY
\begin{frame}{Summary}
\begin{itemize}\setlength{\itemsep}{0.7em}
  \item \code{dl\_complexity} exposes classical, MDL, Bayesian, prequential, and
        compression-based estimators through one transparent interface.
  \item The benchmark relates their native outputs to held-out NLL, the generalization
        gap, computational cost, and rank stability on controlled models and data.
  \item The first result is an end-to-end MLP PoC; four algorithm branches then proceed
        in parallel and undergo cross-review.
  \item Every result retains its target, assumptions, units, components, and warnings;
        common software does not erase theoretical differences.
\end{itemize}

\take{Next technical step: finalize \code{EstimationContext} and
\code{ComplexityResult}, then implement the BIC, two-part MDL, and prequential PoC.}
\end{frame}

% ============================================================ LITERATURE
\begin{frame}[t,shrink=5]{References}
\scriptsize
\begin{thebibliography}{9}\setlength{\itemsep}{0.28em}
\bibitem[Classical criteria]{classical}
H.~Akaike (1974), G.~Schwarz (1978), and E.~Hannan \& B.~Quinn (1979).
Primary definitions of AIC, BIC, and HQIC.

\bibitem[Watanabe(2010, 2013)]{watanabe}
S.~Watanabe. \emph{Asymptotic Equivalence of Bayes Cross Validation and Widely
Applicable Information Criterion in Singular Learning Theory} (JMLR, 2010);
\emph{A Widely Applicable Bayesian Information Criterion} (JMLR, 2013).

\bibitem[Grunwald(2004)]{grunwald}
P.~Grunwald. \emph{A Tutorial Introduction to the Minimum Description Length Principle}.
arXiv:math/0406077, 2004.

\bibitem[Graves(2011)]{graves}
A.~Graves. \emph{Practical Variational Inference for Neural Networks}. NeurIPS, 2011.

\bibitem[Ritter et al.(2018)]{ritter}
H.~Ritter, A.~Botev, D.~Barber. \emph{A Scalable Laplace Approximation for Neural Networks}.
ICLR, 2018.

\bibitem[Bornschein et al.(2022)]{bornschein}
J.~Bornschein, Y.~Li, M.~Hutter. \emph{Sequential Learning of Neural Networks for
Prequential MDL}. arXiv:2210.07931, 2022.

\bibitem[Arora et al.(2018)]{arora}
S.~Arora, R.~Ge, B.~Neyshabur, Y.~Zhang. \emph{Stronger Generalization Bounds for Deep
Nets via a Compression Approach}. ICML, 2018.
\end{thebibliography}
\end{frame}

\end{document}
